%_______________
\end{multicols}



%_______________
\eocesolch{Inferensi untuk data numerik}



%_______________
\begin{multicols}{2}

% 1

\eocesol{(a)~$df=6-1=5$, $t_{5}^{\star} = 2.02$ (kolom dengan jumlah dua ekor 0.10,
baris dengan $df=5$).
(b)~$df=21-1=20$, $t_{20}^{\star} = 2.53$ (kolom dengan jumlah dua ekor 0.02,
baris dengan $df=20$).
(c)~$df=28$, $t_{28}^{\star} = 2.05$.
(d)~$df=11$, $t_{11}^{\star} = 3.11$.}

% 3

\eocesol{(a)~0.085, jangan tolak $H_0$.
(b)~0.003, tolak $H_0$.
(c)~0.438, jangan tolak $H_0$.
(d)~0.042, tolak $H_0$.}

% 5

\eocesol{Rata-rata adalah titik tengah: $\bar{x} = 20$. Tentukan margin galat:
$ME = 1.015$, lalu gunakan $t^{\star}_{35} = 2.03$ dan $SE=s/\sqrt{n}$ dalam
rumus margin galat untuk memperoleh $s = 3$.\\[6mm]}

% 7

\eocesol{(a)~$H_0$: $\mu = 8$ (Warga New York tidur rata-rata 8 jam per malam.)
$H_A$: $\mu \neq 8$ (Warga New York tidur kurang atau lebih dari 8 jam per
malam secara rata-rata.)
(b)~Independensi: Sampelnya acak.
Nilai min/maks menunjukkan tidak ada pencilan yang mengkhawatirkan.
$T = -1.75$. $df=25-1=24$.
(c)~ nilai-p $= 0.093$.
Jika rata-rata populasi sebenarnya untuk
lama tidur warga New York per malam adalah 8 jam,
peluang memperoleh
sampel acak berisi 25 warga New York dengan rata-rata
lama tidur 7.73 jam
per malam atau kurang (atau 8.27 jam atau lebih) adalah 0.093.
(d)~Karena nilai-p $>$ 0.05, jangan tolak $H_0$.
Data tidak memberikan bukti kuat bahwa
warga New York rata-rata tidur lebih atau kurang dari 8 jam
per malam.
(e)~Tidak, karena nilai-p lebih kecil daripada $1 - 0.90 = 0.10$.}

\end{multicols}
\newpage
\begin{multicols}{2}

% 9

\eocesol{$T$ adalah -2.09 atau 2.09.
Kemudian $\bar{x}$ adalah salah satu dari nilai berikut:
\begin{align*}
-2.09 &= \frac{\bar{x} - 60}{\frac{8}{\sqrt{20}}} \ \rightarrow \ \bar{x} = 56.26 \\
2.09 &= \frac{\bar{x} - 60}{\frac{8}{\sqrt{20}}} \ \rightarrow \ \bar{x} = 63.74
\end{align*}}

% 11

\eocesol{(a)~Kita akan melakukan uji $t$ 1-sampel.
    $H_0$: $\mu = 5$. $H_A$: $\mu \neq 5$.
    Kita gunakan $\alpha = 0.05$.
    Ini merupakan sampel acak, sehingga pengamatannya independen.
    Untuk melanjutkan, kita asumsikan distribusi lama mengikuti les
    piano kira-kira normal.
    $SE = 2.2 / \sqrt{20} = 0.4919$.
    Statistik ujinya adalah $T = (4.6 - 5) / SE = -0.81$.
    $df = 20 - 1 = 19$.
    Luas satu ekor kira-kira 0.21, sehingga
    nilai-p kira-kira 0.42, yang lebih besar daripada
    $\alpha = 0.05$, dan kita tidak menolak $H_0$.
    Artinya, kita tidak memiliki bukti yang cukup kuat
    untuk menolak anggapan bahwa rata-ratanya 5 tahun. \\
(b)~Dengan $SE = 0.4919$ dan $t_{df = 19}^{\star} = 2.093$,
    interval kepercayaannya adalah (3.57, 5.63).
    Kita yakin 95\% bahwa rata-rata banyaknya
    tahun seorang anak mengikuti les piano di kota ini berada antara
    3.57 dan 5.63 tahun. \\
(c)~Keduanya konsisten, karena kita tidak menolak hipotesis nol
    dan nilai nol 5 berada dalam interval $t$.}

% 13

\eocesol{Jika sampelnya besar, margin galat akan kira-kira sebesar
$1.96 \times 100 / \sqrt{n}$. Kita ingin nilai ini kurang dari 10, sehingga
diperoleh $n \geq 384.16$; artinya, kita memerlukan ukuran sampel sedikitnya 385 (bulatkan
ke atas untuk perhitungan ukuran sampel!).}
